\section{The Decision Plane: Engines and Algorithms}\label{ch:impl_decision}

The earlier part of this chapter showed how the data plane watches traffic and
how the control plane carries decisions back to the load balancer. This part
opens up the engines themselves: the algorithms that turn raw telemetry into the
three decisions SmartLoad makes (which backend is unhealthy, how much demand is
coming, and where each request should go) plus the forecast-driven autoscaler
that turns a demand prediction into a pool-size change. These are the parts most
readers think of as ``the intelligence'' of the system. They are also where the
engineering choices show most clearly. Every learned engine in SmartLoad ships
next to a deterministic baseline, and in this release the deterministic baseline
is the chosen default in two of the three services. This section explains why
that choice was made and how the architecture turns it into a safe, reversible
decision that the operator controls.

\subsection{The pluggable engine pattern}\label{imb:pluggable}

Every decision-plane service is built on the same skeleton, which we call the
\emph{plugin-host} pattern. A service is a thin run loop that reads telemetry on
a fixed cadence, hands the parsed observation to a swappable \emph{engine}, takes
the engine's structured verdict, and publishes it onto the control bus as a
versioned envelope. The engine is pure: it takes a dataclass of inputs and
returns a dataclass of outputs, and it does no input or output of its own, with
no database query, no Redis publish, no environment-variable read, and no
logging. Every side effect lives in the run loop. This split is enforced
strictly: an engine that opens a socket or builds an envelope is rejected at
review, because the moment a model author can reach the bus directly, the safety
guarantees the run loop provides stop living in one auditable place.

The pattern gives each service three kinds of engine:

\begin{itemize}
  \item A \textbf{deterministic baseline}, a rule with no trained artefact,
        no external dependency, and a constructor that cannot raise. It is
        always present and always correct, if simple.
  \item One or more \textbf{trained engines}, a machine-learning model loaded
        from a serialised artefact at start-up, offering sharper discrimination
        when its training distribution matches production.
  \item An \textbf{automatic fallback} contract that binds the two together: the
        run loop tries to build the requested engine, and if the artefact is
        missing, corrupt, or shape-mismatched, it catches the exception and
        installs the baseline instead. A service is therefore never in a
        ``no engine available'' state.
\end{itemize}

This fallback contract is identical across the three services, because it lives
in one shared helper, \texttt{runloop.bootstrap\_engine()}.
\Cref{fig:imb:engine_bootstrap} traces the decision it makes at start-up. When
the requested engine is a trained model, the helper loads it only if its artefact
is present and valid; if the artefact is missing, corrupt, or shape-mismatched,
the helper catches the failure and installs the deterministic baseline instead.
The baseline has no artefact and a constructor that cannot raise, so it cannot
fail to load. The one case the helper deliberately does not swallow is a baseline
that itself fails to construct, which can only be a deployment bug and is allowed
to surface at start-up. The net effect is that a missing or corrupt model
degrades the service to its baseline rather than to an outage, and the
\texttt{/health} endpoint reports which engine is actually running so the
operator can see when a fallback has occurred.

\begin{figure}[ht]
  \centering
  \includegraphics[width=\linewidth]{engine_bootstrap}
  \caption{Engine bootstrap in \texttt{runloop.bootstrap\_engine()}, the same path
  in every decision-plane service. A trained model is installed when its artefact
  is present and valid; otherwise the deterministic baseline takes over. Because
  the baseline has no artefact and cannot fail to construct, a missing or corrupt
  model degrades the service to the baseline rather than to an outage, and the
  resolved engine is surfaced on \texttt{/health}.}
  \label{fig:imb:engine_bootstrap}
\end{figure}

The operator picks the engine through a single environment variable
(\texttt{ANOMALY\_ENGINE}, \texttt{FORECAST\_ENGINE}, \texttt{RL\_POLICY}), and
the run loop owns the envelope, the publication-gating logic, the policy
subscription, and the fallback decision. Adding a new engine is one folder and
one factory branch, with zero edits to the run loop. This is what makes the
design posture of the next sections work in practice: choosing which engine leads
in production is a one-line configuration change, not a code rollback, and the
baseline that takes over is the same code path that has run since the first
sprint.

The run loop itself follows a fixed cycle, shown in \cref{fig:imb:runloop_seq}.
On each tick it first drains any pending policy update so a configuration change
takes effect at the next cycle, then queries telemetry, builds the engine's input
features, runs the engine, and finally publishes a decision envelope onto the
control bus, but only when a should-publish gate agrees. The gate is what
suppresses redundant or unsafe output: a verdict is dropped, and a counter
incremented, when the service is in safe mode or when the result carries no new
information. The tick interval and the telemetry query differ from one service to
the next, while the control flow itself is shared, which is what lets the safety
behaviour live in one reviewed place rather than being re-implemented per engine.

\begin{figure}[p]
  \centering
  \includegraphics[width=\linewidth,height=0.9\textheight,keepaspectratio]{runloop_seq}
  \caption{The per-service run-loop cycle. On each tick the loop reads any pending
  policy update, queries telemetry, builds features, runs the engine, and
  publishes a decision envelope only when the should-publish gate passes;
  otherwise the verdict is dropped and a counter incremented. The tick interval
  and the query differ per service (anomaly, forecast, routing), while the control
  flow is shared.}
  \label{fig:imb:runloop_seq}
\end{figure}

\subsection{Why these algorithms, and not others}\label{imb:algo_choices}

The algorithm surface of the decision plane concentrates into five choices, one
per learned engine plus the framing of the routing model. Each was a real
trade-off between a simpler option that ships immediately and a heavier option
that promises more discrimination at a higher integration cost.
\Cref{tab:imb:algo_choices} records the five decisions with the alternative that
was set aside and the reason. The pattern is consistent: SmartLoad adopts the
lightest implementation that clears the relevant acceptance gate, and keeps the
heavier candidate as a documented upgrade rather than a blocker.

\begin{table}[ht]
  \centering
  \caption{Algorithm-selection decisions across the decision plane.}
  \label{tab:imb:algo_choices}
  \small
  \begin{tabularx}{\linewidth}{@{}p{2.5cm}p{2.9cm}p{2.7cm}Y@{}}
    \toprule
    \headrow
    \textbf{Decision} & \textbf{Chosen} & \textbf{Rejected alternative} & \textbf{Reason} \\
    \midrule
    Anomaly baseline & Threshold rule & Isolation Forest as the day-one baseline
      & The threshold rule ships with no model or training data, while the forest
        needs an artefact. Making the rule the baseline keeps the forest's own
        fallback path exercised on every cycle. \\
    \addlinespace
    Anomaly upgrade & Isolation Forest & LSTM autoencoder
      & The forest is a single package, trains quickly, and serialises to a small
        pickle. The autoencoder catches contextual anomalies the forest misses,
        but its cost (GPU at inference, a deeper training pipeline, a per-workload
        training set) is not warranted at an $F_1>0.80$ gate. It stays a deferred
        upgrade. \\
    \addlinespace
    Forecast baseline & Moving average & A per-service ``no forecast'' state
      & A windowed mean is always available, which removes the
        ``no forecast'' failure mode and gives the autoscaler a usable signal
        under every condition. \\
    \addlinespace
    Forecast model & ARIMA$(3,0,1)$ & Prophet, DeepAR
      & The five-minute horizon does not yet demand the extra modelling capacity.
        ARIMA's footprint is well understood and the artefact ships trivially;
        the heavier models are deferred. \\
    \addlinespace
    Routing framing & Offline contextual bandit on logged traces & Online MDP with a learned environment model
      & Reproducibility, safety, and training cost favour the offline framing. The
        accepted cost is that the policy does not reason about its own routing
        impact, an axis the production feedback loops cover instead
        (\cref{imb:ppo_mdp}). \\
    \bottomrule
  \end{tabularx}
\end{table}

\section{Data Sources and Tooling}\label{imb:data_tooling}

The learned engines are only as good as the data they were trained on and the
libraries that implement them. This section records both, because the training
corpus is the biggest single factor in how well each engine carries over to
SmartLoad's mixed test workload, a theme that returns through the rest of the
chapter.

\subsection{Datasets}\label{imb:datasets}

SmartLoad draws on four public datasets, each chosen for a specific engine and a
specific reason.

\textbf{Google Borg cluster traces} \cite{verma2015borg, tirmazi2020borg}
provide month-long records of task scheduling, resource requests, and machine
utilisation across a production warehouse-scale cluster. We use them as the
reference shape for realistic load-arrival patterns and to ground the
autoscaler's capacity bands against utilisation seen in a real cluster, rather
than against a synthetic curve of our own invention.

\textbf{Alibaba microservice traces} \cite{guo2019alibaba} record per-call
latencies and call-graph topology across a large microservice deployment. These
are the training corpus for the learned routing policy: each trace window is
replayed as a set of per-backend latency snapshots, from which the policy learns
which backend is likely to be fastest in the next interval. One useful property
of these traces, which shapes the routing discussion in \cref{imb:ppo_finding},
is that the latency across backends is fairly \emph{uniform}.

\textbf{The Server Machine Dataset (SMD)} \cite{su2019omnianomaly} is a labelled
multivariate anomaly-detection benchmark of 38 machine-level metric streams with
ground-truth anomaly intervals. It was the first training corpus for the anomaly
engine, chosen because it ships real anomaly labels, a rarity in this space, and
so supports a defensible $F_1$ score on a held-out partition.

The Borg and Alibaba corpora are large and are kept out of version control by
design; the engines consume them through training pipelines rather than at
runtime, and the runtime image carries only the resulting model artefacts.

\subsection{Software tools and libraries}\label{imb:libraries}

\Cref{tab:imb:tooling} lists the main libraries in the decision plane and the
reason for each. The guiding principle is that every learned engine uses a
mature, widely audited implementation of its algorithm rather than a hand-rolled
one: the research contribution of SmartLoad is the integration architecture and
the closed-loop evaluation, not a new numerical kernel.

\begin{table}[ht]
  \centering
  \caption{Decision-plane libraries and their justification.}
  \label{tab:imb:tooling}
  \small
  \begin{tabular}{@{}llp{3.0cm}p{4.4cm}@{}}
    \toprule
    \textbf{Component} & \textbf{Library} & \textbf{Role} & \textbf{Justification} \\
    \midrule
    Anomaly (trained)  & scikit-learn        & Isolation Forest model + scalers
      & Mature, deterministic unsupervised outlier detector; stable pickle format pinned to one version. \\
    Forecast (trained) & statsmodels         & ARIMA estimation + forecasting \cite{box2015timeseries}
      & Reference implementation of classical time-series models with native confidence intervals. \\
    Routing (trained)  & Stable-Baselines3 / sb3-contrib & MaskablePPO policy \cite{schulman2017ppo, raffin2021sb3}
      & Audited PPO implementation with action-masking, removing the need to hand-roll the optimiser. \\
    Routing backend    & PyTorch             & Neural-network forward pass
      & The tensor backend Stable-Baselines3 builds on; one forward pass per routing decision. \\
    Persistence        & psycopg2            & TimescaleDB access in the run loop
      & The canonical PostgreSQL driver; all engine inputs come from parameterised SQL it executes. \\
    Control bus        & redis-py            & Envelope publish / subscribe
      & Lightweight pub/sub the run loop uses to carry verdicts to the data plane. \\
    Observability      & prometheus-client   & Per-process \texttt{/metrics} surface \cite{prometheus}
      & Standard exposition format scraped by Prometheus for engine-cycle and publish counters. \\
    \bottomrule
  \end{tabular}
\end{table}

The deferred-import discipline in the engine factory is worth a sentence here. A
heavy dependency such as \texttt{statsmodels}, \texttt{scikit-learn}, or
\texttt{torch} is imported only inside the factory branch that actually builds
the engine that needs it. So the baseline path never pays the import cost of a
library it does not use, and a service running its baseline does not need the
trained engine's dependencies to be present at all.

\section{Anomaly Detection}\label{imb:anomaly}

The anomaly-detector classifies each backend, every cycle, into one of three
states (\texttt{healthy}, \texttt{degraded}, or \texttt{unhealthy}) from a
short window of per-backend latency and error-rate telemetry. The
\texttt{unhealthy} verdict matters: it is published on the
\texttt{smartload.anomaly} channel and tells the load-balancer sidecar to remove
the backend from the upstream pool until it recovers. The service ships two
engines.

\subsection{Threshold baseline}\label{imb:threshold}

The baseline is a deterministic rule of roughly forty lines. It is the production
default in the shipped configuration, the automatic fallback if a trained engine
cannot load, and the reference signal the trained engine is calibrated against. Its logic, in order of severity, is:

\begin{enumerate}
  \item If the sample count in the window is below a minimum, report
        \texttt{healthy} (not enough data to judge).
  \item If the observed error rate goes over a fixed threshold, report
        \texttt{unhealthy}. Error rate is checked first because it is a harder
        safety signal than latency.
  \item Otherwise, if the current latency goes over a multiple of the backend's
        own rolling-mean latency, report \texttt{degraded}.
  \item Otherwise, report \texttt{healthy}.
\end{enumerate}

Formally, writing $\ell$ for the current latency, $\bar{\ell}$ for the rolling
mean, $e$ for the error rate, $m$ for the latency multiplier (default $3.0$),
and $\tau$ for the error-rate threshold (default $0.05$), the backend is flagged
when

\begin{equation}\label{imb:eq:threshold}
  \texttt{unhealthy} \iff e > \tau,
  \qquad
  \texttt{degraded} \iff \frac{\ell}{\bar{\ell}} > m \;\wedge\; e \le \tau .
\end{equation}

The rule compares each backend against \emph{its own} recent history rather than
against a global latency constant. A backend that is genuinely slow but stable is
therefore not penalised; only a backend that is slow relative to its own
baseline, or that is erroring, gets flagged. This relative formulation is the
same reference signal the trained engine is later calibrated to agree with.

\subsection{Isolation Forest engine}\label{imb:isoforest}

The trained engine is an Isolation Forest \cite{liu2008isolationforest}, an
unsupervised ensemble that isolates outliers by random recursive partitioning:
points that need few splits to isolate are anomalous. We chose it because it
makes no parametric assumption about the shape of the ``normal'' region, trains
quickly, and scores in constant time per point.

\paragraph{Training and validation.} The engine was trained to recognise the
shape of healthy SmartLoad telemetry and reached an $F_1$ of $0.8012$ on a
held-out partition, clearing the project's $>0.80$ acceptance gate. That is a
solid pass of the engine's own validation test, and it produced a trained
detector that could be lined up directly against the threshold rule on
production-shaped inputs.

\paragraph{Calibration in production-shape space.} We fit the model and its
scaler in one shared real-millisecond coordinate system, on a synthetic
\emph{healthy operating region} that encodes what a healthy SmartLoad backend
emits: latency spanning the full operating range, consistent request behaviour
(current latency close to the rolling mean), error rate at or below the
baseline's boundary, and a characteristic variance band. The threshold rule
served as the reference signal, and the decision thresholds were searched
empirically against a bench sweep of $(\text{latency}, \text{error rate})$ cells
spanning the operating range. The trained engine agrees with the threshold
baseline on \textbf{91.4\%} of cells (234 of 256), and every one of the 22
remaining disagreements is in the value-adding direction: the threshold rule says
\texttt{healthy} but the model catches a variance or spike pattern the rule
misses. A good example is the live-stack 400\,ms latency injection, whose
measured feature vector $(\text{max}=406, \text{avg}=162, \text{err}=0,
\text{std}=196)$ is flagged \texttt{unhealthy} by the model on its
variance-to-mean ratio of roughly $1.2$, a signal the threshold rule's latency
ratio of $2.5 < 3$ does not pick up.

\Cref{fig:imb:scoring_pipeline} lays out the full scoring pipeline and makes the
coordinate-system bridge explicit. At training time the model is fitted on the
SMD corpus in its own normalised space, while a second scaler, the
\texttt{production\_scaler}, is fitted on real-millisecond features so that the
two spaces can be reconciled. Both scalers, the model, the feature order, and the
derived thresholds are serialised into a single bundle. At serving time the
engine assembles the raw feature vector, passes it through the
\texttt{production\_scaler} to map live real-millisecond telemetry into the space
the model was trained in, runs the decision function, and classifies the result
against the bundled thresholds. The bridge between the training feature space and
the live serving feature space is the load-bearing detail: without it the model
would be asked to score inputs in a coordinate system it never saw. This is the
same train/serve construct gap that recurs across the learned engines, surfacing
again for the routing policy in \cref{imb:ppo_mdp}, and here the calibration step
is what closes it for the anomaly detector.

\begin{figure}[ht]
  \centering
  \includegraphics[width=\linewidth]{scoring_pipeline}
  \caption{The Isolation Forest scoring pipeline and the coordinate-system bridge
  between the training feature space and the live serving feature space. At
  training time the model is fitted on the SMD corpus in its own normalised space,
  and a separate production scaler is fitted on real-millisecond features; both,
  with the model, feature order, and thresholds, are serialised into one bundle.
  At serving time the engine maps live telemetry through the production scaler into
  the model's coordinate system before scoring and classifying it. The bridge is
  what lets a model trained in one feature space score inputs drawn from another.}
  \label{fig:imb:scoring_pipeline}
\end{figure}

\paragraph{Deployment status.} The deterministic threshold rule is the production
default, chosen for guaranteed stability and because it doubles as the automatic
fallback for the whole service. The Isolation Forest engine is available as an
enhanced detection mode via \texttt{ANOMALY\_ENGINE=isolation\_forest}: it agrees
closely with the threshold rule and, on top of that, catches extra latency
variance and error-rate compounding that the simple rule alone would miss. We
also added a quorum guard in the sidecar that refuses any exclusion which would
empty the pool, which protects both engines equally. This is the pluggable
pattern working as designed: an enhanced detector is available behind a flag, the
deterministic baseline leads in production for safety, and switching between them
is a one-line configuration change.

\Cref{lst:imb:if_score} gives the scoring loop the engine runs each cycle.

\begin{lstlisting}[language=Python, caption={Isolation Forest scoring loop (per backend, per cycle).}, label={lst:imb:if_score}]
def score(features):
    # Data-quality gate: too few samples to judge.
    if features.sample_count < min_sample_count:
        return AnomalyScore(features.backend_id, "healthy", 0.0)

    # Assemble the raw feature vector in the bundle's feature order
    # (real-millisecond latency, error fraction, variance).
    x_raw = assemble(features, feature_order)

    # Map into the model's coordinate system. After calibration the
    # scaler and the model share one real-ms space, so this is identity-like.
    x = production_scaler.transform([x_raw])

    # Lower decision_function => more anomalous.
    raw = model.decision_function(x)[0]

    if raw > thresholds.healthy_above:
        return AnomalyScore(features.backend_id, "healthy", 0.0)
    if raw >= thresholds.unhealthy_below:
        return AnomalyScore(features.backend_id, "degraded", 0.5)

    severity = min(1.0, abs(raw - thresholds.unhealthy_below)
                        / thresholds.unhealthy_score_scale)
    return AnomalyScore(features.backend_id, "unhealthy", severity)
\end{lstlisting}

\section{Demand Forecasting}\label{imb:forecast}

The forecasting service projects the recent request-rate trend forward by a fixed
horizon (five minutes by default) and publishes a predicted requests-per-second
value with a confidence band. The autoscaler consumes that prediction. As with
anomaly detection, the service ships a deterministic baseline and a trained
engine. The moving-average baseline is the reliable production default and is the
predictor that drives the demonstrated autoscaling loop, while the ARIMA engine
is an optional advanced predictor we keep tuning.

\subsection{Moving-average baseline}\label{imb:movingavg}

The baseline is a windowed mean \cite{hyndman2021forecasting}. It takes the last
$k$ one-minute request-rate buckets, returns their mean as the point prediction,
and returns one sample standard deviation either side as a flapping-suppression
band rather than a true probabilistic interval. With $r_1, \dots, r_k$ the recent
per-bucket rates, the prediction is simply

\begin{equation}\label{imb:eq:movingavg}
  \hat{r} = \frac{1}{k}\sum_{i=1}^{k} r_i,
  \qquad
  [\,\hat{r}-s,\ \hat{r}+s\,], \quad
  s = \sqrt{\tfrac{1}{k-1}\textstyle\sum_{i=1}^{k}(r_i-\hat{r})^2}.
\end{equation}

This is deliberately simple. Its strength, and the reason it ships as a
first-class engine rather than a placeholder, is that it is always available and
never fails. That removes the ``no forecast available'' state and gives the
autoscaler a usable signal under every condition.

\subsection{ARIMA(3,0,1)}\label{imb:arima}

The trained engine is an ARIMA$(3,0,1)$ model \cite{box2015timeseries}: three
autoregressive terms, no differencing, one moving-average term, fit offline on
historical request-rate data and serialised as a model artefact. At serving time
the engine appends the recent window to the fitted model's state without
refitting the coefficients, an $O(\text{window})$ update, and forecasts one
step ahead, clamping the prediction and its confidence bounds to be non-negative,
since a request rate cannot fall below zero.

The model's accuracy is measured by the mean absolute percentage error over the
held-out partition. With $y_t$ the actual rate and $\hat{y}_t$ the prediction
over $n$ points,

\begin{equation}\label{imb:eq:mape}
  \mathrm{MAPE} = \frac{100\%}{n}\sum_{t=1}^{n}
  \left| \frac{y_t - \hat{y}_t}{y_t} \right| .
\end{equation}

\paragraph{Status.} ARIMA is wired and serving-ready behind
\texttt{FORECAST\_ENGINE=arima}. The current artefact measures about
$25.0\%$ test MAPE against a $<20\%$ acceptance gate, so it does not yet clear
its own validation and remains under continued tuning through hyperparameter
search and an extended training window. The moving-average engine stays the
production default because it is the dependable signal that drives the autoscaler
today: the forecasting path works end-to-end and steers the autoscaler
correctly, and ARIMA sits alongside it as the optional advanced predictor under
development.

\section{Learned Routing with PPO}\label{imb:ppo}

The routing engine (the \texttt{rl-engine}) is where SmartLoad's most ambitious
learned component lives. The engine ranks the backend pool each cycle and
publishes a routing recommendation; when, and only when, every safety gate
agrees, the sidecar turns that ranking into NGINX upstream weights. By default
the engine runs in \emph{shadow} mode: it computes and publishes its
recommendation so the decision is observable, but the sidecar does not act on it.
Around the learned policy we ship three proven deterministic baselines, all
behind one pluggable interface.

\subsection{Observation, action, and action masking}\label{imb:ppo_spec}

The learned policy observes, for each backend, its recent latency, a queue-depth
proxy (the windowed request count, since the schema carries no true connection
queue), and an anomaly-derived health flag. The three features are normalised
into a fixed-length tensor, and the policy emits one chosen backend slot per
step. \Cref{tab:imb:ppo_spec} gives the exact specification. The same
\texttt{build\_observation} and \texttt{build\_action\_mask} functions run at
training time and at serving time, which guarantees train/serve parity: the
serving plugin restores the normalisation constants from the artefact metadata so
the policy sees inputs scaled exactly as it was trained on.

\begin{table}[ht]
  \centering
  \caption{PPO observation and action specification (\texttt{rl-engine}).}
  \label{tab:imb:ppo_spec}
  \small
  \begin{tabularx}{\linewidth}{@{}p{2.5cm}p{3.5cm}Y@{}}
    \toprule
    \headrow
    \textbf{Element} & \textbf{Definition} & \textbf{Notes} \\
    \midrule
    Observation & $\mathrm{Box}(N_{\max}\!\times\!3)$, per backend the triple
      latency, request count, and health flag
      & Each feature is normalised: latency by $2000$\,ms (Alibaba p99); request
        count by $200$ (a typical 30\,s window); health flag is $1.0$ if eligible,
        else $0.0$. Fixed slot order, sorted by backend identifier. \\
    \addlinespace
    Action & $\mathrm{Discrete}(N_{\max})$
      & The index of the single backend the policy selects each step;
        $N_{\max}=5$ matches the production pool. \\
    \addlinespace
    Action masking & Invalid-backend logits set to $-\infty$ before the softmax
      & Eligible slots (\texttt{healthy} or \texttt{degraded}) stay; unhealthy,
        unknown, and empty slots are masked, so the policy assigns them zero
        probability and can never propose an unhealthy backend
        \cite{schulman2017ppo, raffin2021sb3}. \\
    \bottomrule
  \end{tabularx}
\end{table}

The policy is optimised with MaskablePPO \cite{schulman2017ppo, raffin2021sb3},
the action-masked variant of Proximal Policy Optimisation, so that ineligible
backends are removed from the action distribution rather than merely penalised.
Masking is the load-bearing safety property of the learned routing path: it is a
hard structural guarantee that the policy cannot select a backend the anomaly
detector has flagged, independent of whatever the network weights have learned.

\subsection{The reward function}\label{imb:ppo_reward}

The per-step reward is recorded in one canonical file and never duplicated in
serving code. Writing $s'_a$ for the chosen backend as observed in the next state,
the reward is

\begin{equation}\label{imb:eq:reward}
  r(s, a, s') =
  -\,\frac{\ell(s'_a)}{\texttt{latency\_scale}}
  \;-\; \lambda \cdot \mathrm{imbalance}(s')
  \;-\; 10.0 \cdot \mathbb{1}[\,s'_a.\texttt{health} = \texttt{unhealthy}\,],
\end{equation}

with $\texttt{latency\_scale}=2000$\,ms and $\lambda=0.1$. \Cref{tab:imb:reward}
explains each term and why it is present. \Cref{imb:eq:reward} evaluates the
chosen backend on the \emph{next} state, which gives the policy correct credit
assignment: it is rewarded for picking a backend that turns out to be fast in the
following window, not for the state it already saw.

\begin{table}[ht]
  \centering
  \caption{Terms of the PPO reward function (\cref{imb:eq:reward}).}
  \label{tab:imb:reward}
  \small
  \begin{tabularx}{\linewidth}{@{}p{3.7cm}Y@{}}
    \toprule
    \headrow
    \textbf{Term} & \textbf{Role and rationale} \\
    \midrule
    Latency penalty
      $-\,\ell(s'_a)/\texttt{latency\_scale}$
      & The primary objective. Scaling by the same $2000$\,ms used for the
        observation keeps the latency signal on a stable order of magnitude
        across episodes. \\
    \addlinespace
    Imbalance regulariser
      $-\,\lambda\cdot\mathrm{imbalance}(s')$
      & $\mathrm{imbalance}(s')=\mathrm{std}(\text{queue depths})/(\mathrm{mean}+\varepsilon)$,
        $\varepsilon=1.0$. Because the simulator does not model how a routing
        choice perturbs queue depths, this is a regulariser on the historical
        spread, not a learned consequence. It discourages picks biased toward
        load outliers and smooths the training signal. $\lambda=0.1$ is the
        largest weight at which removing the term did not change held-out metrics
        in the canary. \\
    \addlinespace
    Health penalty
      $-\,10.0\cdot\mathbb{1}[\text{unhealthy}]$
      & A defence-in-depth hard penalty. In practice the action mask already
        prevents an unhealthy pick, so this term only fires if the mask is ever
        misconfigured, where it supplies a strong negative signal. \\
    \bottomrule
  \end{tabularx}
\end{table}

\subsection{Why a contextual bandit, not a full MDP}\label{imb:ppo_mdp}

It helps to be precise about what the trained artefact is. It is a
\emph{contextual bandit} rather than a full reinforcement-learning agent in the
Markov-decision-process sense \cite{sutton2018rl}. The offline training simulator
replays trace windows independently of the action the policy takes, so
\texttt{step(action)} does not modify the next state. The policy therefore learns
to \emph{predict} which backend will be fast next rather than to model the
downstream effect of its own routing on how the system evolves.

This framing was a deliberate engineering decision rather than a limitation
discovered late, and the trade-off is recorded in \cref{tab:imb:algo_choices}.
An online MDP with a learned environment model would let the policy reason about
its own routing impact, but at a cost in reproducibility, safety, and training
expense that the project judged unwarranted at this stage. The offline bandit
trains on logged production traces under a clean temporal split, which keeps the
held-out evaluation honest and the result reproducible. The closed-loop
``consequence'' axis is supplied instead by the deterministic safety machinery
around the policy: NGINX's own failure detection, the anomaly detector's
exclusions, and the autoscaler's reactivity. This framing is also what makes the
shadow-mode evaluation of \cref{imb:ppo_finding} interpretable, and its
consequence for the policy's measured behaviour is taken up in
\cref{ch:discussion}.

\subsection{Serving: argmax-dominant weighting}\label{imb:ppo_serving}

The policy was trained to pick one backend per step, but the load balancer needs
a weight \emph{distribution} over the pool. Two natural mappings from a
$\mathrm{Discrete}(N)$ policy to upstream weights exist, and the serving path
rejects one and adopts the other; \cref{tab:imb:weight_map} records the choice.

\begin{table}[ht]
  \centering
  \caption{Mapping the discrete policy output to NGINX upstream weights.}
  \label{tab:imb:weight_map}
  \small
  \begin{tabularx}{\linewidth}{@{}p{4.2cm}p{1.9cm}Y@{}}
    \toprule
    \headrow
    \textbf{Mapping} & \textbf{Decision} & \textbf{Reason} \\
    \midrule
    Softmax of the backend logits into a weight distribution & Rejected
      & A $\mathrm{Discrete}(N)$ policy is trained to \emph{pick one} backend, not
        to distribute calibrated mass. The logit distribution was never an
        optimisation target, so serving from it would route traffic against the
        policy's training objective. \\
    \addlinespace
    Argmax-dominant: the chosen backend gets \texttt{\_DOMINANT\_WEIGHT}, the
    remainder split evenly across the other eligibles & Adopted
      & Matches what the policy was optimised for. A single dominant weight of
        $0.7$ emphasises the policy's choice while the $0.3$ floor keeps NGINX
        health probing live on the rest of the pool. \\
    \bottomrule
  \end{tabularx}
\end{table}

Reusing the policy's raw output distribution as weights would not match what it
was optimised for, so the serving path instead assigns a dominant weight (default
$0.7$) to the chosen (argmax) backend and splits the remaining mass evenly across
the other eligible backends. The small floor is deliberate: a backend with zero
weight never receives traffic and so never reveals whether it has recovered,
whereas a small share keeps NGINX's health probing live on every eligible
backend. \Cref{lst:imb:ppo_rank} gives the ranking loop.

\begin{lstlisting}[language=Python, caption={PPO serving: one forward pass, then argmax-dominant weighting.}, label={lst:imb:ppo_rank}]
def act(state):
    eligible = [b for b in state if is_eligible(b.health)]
    if not eligible:
        return routing_fallback(state)   # canonical no-eligible path

    obs  = build_observation(state, N_MAX_BACKENDS, norm_params)
    mask = build_action_mask(state, N_MAX_BACKENDS)

    logits        = forward_pass(obs)              # single inference
    masked_logits = where(mask, logits, -inf)      # ineligible -> -inf
    chosen_slot   = argmax(masked_logits)

    n_other = len(eligible) - 1
    floor   = (1.0 - DOMINANT_WEIGHT) / n_other if n_other > 0 else 0.0
    rankings = [
        Ranking(b.backend_id,
                DOMINANT_WEIGHT if b is chosen else floor)
        for b in eligible
    ]
    rankings.sort(key=lambda r: r.score, reverse=True)

    mode = "active" if operating_mode == "hybrid" else "shadow"
    return RoutingAction(mode=mode, rankings=rankings)
\end{lstlisting}

\subsection{The three classical baselines}\label{imb:ppo_baselines}

Alongside the trained policy, the engine ships three deterministic baselines,
each selectable as a policy and each a real comparator rather than a stub:

\begin{itemize}
  \item \textbf{\texttt{random\_shadow}}: uniform-random ranking, always in
        shadow mode. A sanity baseline for the trained policy to clear, and the
        fallback installed when \texttt{policy.zip} cannot be loaded.
  \item \textbf{\texttt{round\_robin}}: stable rotation across eligible
        backends, using a backend-identifier pointer rather than a modular index
        so the cycle does not jump when an anomaly changes the eligible set
        mid-rotation. This is the direct analogue of NGINX's own default.
  \item \textbf{\texttt{least\_connections}}: ranks by ascending queue-depth
        proxy, sending traffic to the backend with the fewest in-flight requests
        in the window.
\end{itemize}

\subsection{Shadow-mode evaluation}\label{imb:ppo_finding}

The trained PPO policy was evaluated in shadow mode against the deterministic
baselines over a held-out seed bank of twenty episodes. The headline result is
that \emph{the trained policy ties NGINX round-robin}, the strongest baseline: the
two share a mean reward of $-0.0056$ and are joint-best, while the policy beats
\texttt{least\_connections} by $4.8\times10^{-2}$ in mean reward. The
SLO-violation rate is $0$ for \emph{every} policy, because the trace backend
latencies (around $9.7$\,ms) sit uniformly below the $200$\,ms objective, so no
routing choice can violate it. Round-robin is a strong, industry-standard
baseline, so a byte-for-byte tie means the learned policy is at least as effective
and as safe as the established approach, and the shadow pin guarantees it can
never do worse in production. Two structural points explain the tie precisely, and
the deeper reading of what it implies for the routing question is taken up in
\cref{ch:discussion}.

\begin{enumerate}
  \item \textbf{Training distribution.} The policy was trained on the Alibaba
        microservice traces, whose per-backend latency is fairly uniform. On a
        homogeneous pool the optimal routing policy \emph{is} uniform sampling,
        which is exactly round-robin: there is no faster backend to prefer. The
        policy learned the correct optimum for its training distribution; the data
        carried no latency signal to discriminate on. SmartLoad's test workload is
        deliberately varied, so retraining on a heterogeneous-latency distribution
        is the named binding constraint for any advantage on varied traffic.
  \item \textbf{Open-loop simulator.} As described in \cref{imb:ppo_mdp}, the
        simulator replays windows independently of the action, so the policy
        optimised a one-step prediction objective. Moving to a closed-loop
        simulator that lets a routing choice perturb the next state is the natural
        next step for the policy to learn the downstream effect of its own
        routing.
\end{enumerate}

The interpretation here is engineering, not failure: the policy converged to
round-robin equivalence because uniform sampling is optimal on a homogeneous
training distribution, and the evaluation that surfaced this is a clean,
temporally split, honestly reported one. The serving mechanism, the safety gates,
and the bench harness are all in place and verified, so the engine already
produces an observable, auditable decision stream that costs the production
routing path nothing. Routing and scaling are separate paths, so the shadow
policy adds no load to the demonstrated autoscaling loop while its training is
extended.

\paragraph{The shadow safety pin.} This posture is underwritten by a hard safety
default: \texttt{RL\_MODE=shadow}. Even with the run loop enabled and a trained
policy loaded, the sidecar refuses to apply any recommendation whose mode is not
\texttt{active}, and the run loop will not emit \texttt{active} unless the
operator has explicitly opted in \emph{and} the operating policy is not in
safe-mode. \Cref{fig:imb:mode_composition} shows how these conditions combine:
three independent gates must all agree before the published mode is
\texttt{active}, so the same intent, namely ``do not move traffic the operator did
not approve'', is enforced redundantly, with the sidecar as the single
authoritative safety boundary. The routing engine therefore only moves production
traffic when the operator has approved it; otherwise it just produces telemetry.
The asymmetry is deliberate: the policy is what learns, while the conditions that
permit it to move traffic are fixed, so routing can improve over time without ever
loosening the safety default.

\begin{figure}[ht]
  \centering
  \includegraphics[width=0.62\linewidth]{mode_composition}
  \caption{How the published routing mode is composed. Three independent gates
  must all agree before any routing recommendation reaches the load balancer as
  \texttt{active}: the operating policy must not be in safe mode (the kill
  switch), the \texttt{RL\_MODE} environment pin must be \texttt{active} (the
  operator opt-in), and the policy itself must return an \texttt{active} action.
  If any gate declines, the published mode is \texttt{shadow} and the sidecar
  ignores the weights. Shadow is the default, and no single component can escalate
  to active on its own.}
  \label{fig:imb:mode_composition}
\end{figure}

\subsection{The load-balancer sidecar}\label{imb:sidecar}

The component that actually applies a routing decision is the load-balancer
sidecar. It subscribes to the routing, anomaly, and policy channels, translates a
recommendation into NGINX upstream directives, and reloads the running NGINX
without touching the request path it serves. Several small design decisions make
that bridge safe to operate, and \cref{tab:imb:sidecar} records them. The theme
is the same one that runs through the rest of the decision plane: each choice
favours an idempotent, recoverable mechanism over a clever one.

\begin{table}[ht]
  \centering
  \caption{Load-balancer sidecar design decisions.}
  \label{tab:imb:sidecar}
  \small
  \begin{tabularx}{\linewidth}{@{}p{4.0cm}Y@{}}
    \toprule
    \headrow
    \textbf{Decision} & \textbf{Rationale} \\
    \midrule
    Rewrite an include file, not \texttt{nginx.conf}
      & Only one upstream file changes per cycle, so the rest of the operator's
        NGINX configuration is never rewritten or at risk. \\
    \addlinespace
    Reload via \texttt{docker exec}, not an HTTP API
      & Open-source NGINX has no reload API. The autoscaler already drives Docker
        the same way, so the mechanism is an established one in the stack. \\
    \addlinespace
    Atomic write (\texttt{os.replace})
      & The temp-file-then-rename is atomic on POSIX, so NGINX always reads a
        complete file on the reload that follows; it never sees a half-written
        upstream block. \\
    \addlinespace
    IP-to-container-name mapping
      & NGINX logs resolved \texttt{IP:port}, so rankings carry IP-based backend
        identifiers. A live registry maps them back to container names and
        refreshes on any unmapped lookup, so new backends are picked up without a
        restart. \\
    \addlinespace
    Weight floor at $1$
      & NGINX rejects \texttt{weight=0}, so any zero score is promoted to $1$.
        This keeps every eligible backend in the rotation and preserves health
        probing on it. \\
    \addlinespace
    Confidence-threshold gate
      & When the top recommendation's score falls below the policy's
        \texttt{rl\_confidence\_threshold}, the recommendation is rejected and the
        previous upstream state stands, so a low-confidence cycle never moves
        traffic. \\
    \bottomrule
  \end{tabularx}
\end{table}

\section{Forecast-Driven Autoscaling}\label{imb:autoscaler}

The autoscaler closes the loop from prediction to capacity. It subscribes to the
forecast channel and to the policy channel, turns a predicted request rate into a
target backend count, and actuates the pool through the Docker SDK. Importantly,
anomaly signals do \emph{not} drive scaling. Capacity decisions are
forecast-driven by design, while anomaly verdicts drive routing exclusion.
Keeping these channels separate stops a brief health blip from triggering a
capacity change.

\subsection{The scaling decision}\label{imb:autoscaler_decision}

The core decision is a capacity comparison with built-in hysteresis. With
$\hat{r}$ the predicted rate, $n$ the current backend count, and $c$ the
per-instance capacity in requests per second drawn from policy, the autoscaler
scales out when demand goes over current capacity and scales in only when
shedding one backend would \emph{still} leave capacity above demand:

\begin{equation}\label{imb:eq:scale}
  \text{action} =
  \begin{cases}
    \texttt{scale\_out} & \text{if } \hat{r} > n\,c \\[2pt]
    \texttt{scale\_in}  & \text{if } \hat{r} < (n-1)\,c \\[2pt]
    \texttt{noop}       & \text{otherwise (within the band } [(n-1)c,\ nc]).
  \end{cases}
\end{equation}

The conceptual target count is $\lceil \hat{r} / c \rceil$; the engine moves one
backend at a time toward it. The dead band $[(n-1)c,\ nc]$ is what avoids
oscillation: a prediction that hovers around a backend boundary does not make the
pool flap, because scaling in requires demand to fall a full instance's worth
below current capacity. Every decision also respects a cooldown timer (no action
within \texttt{autoscaler\_cooldown\_seconds} of the last), and the resulting
count is clamped to the policy's \texttt{min\_backends} and \texttt{max\_backends}
bounds.

When the forecast stream goes stale, with no fresh prediction inside a staleness
window, the autoscaler falls back to a \emph{reactive} mode, computing the
same decision against an observed request rate read directly from telemetry
rather than against a forecast. The audit reason field records which path drove
each decision, so a reader of the scaling history can always tell a proactive
scale from a reactive one.

\subsection{Actuation and audit}\label{imb:autoscaler_actuation}

Actuation goes through one of two lifecycle pairs in the cluster client. The
default pair, \texttt{start()}/\texttt{stop()}, toggles compose-provisioned
containers and is used by the routing benchmark. The dynamic pair,
\texttt{provision()}/\texttt{decommission()}, creates and destroys containers
through the Docker SDK and is gated behind a feature flag for the adaptive
benchmark. Each scaling action publishes a \texttt{ScalingEvent} on the
\texttt{smartload.scale} channel, which the sidecar consumes to rewrite the
NGINX upstream to match the live pool, and writes a row to the
\texttt{scaling\_events} audit table recording the action, the resulting count,
the mechanism, and the human-readable reason.

\Cref{fig:imb:autoscaler} shows the decision flow.

\begin{figure}[ht]
  \centering
  \begin{tikzpicture}[
    node distance=8mm and 14mm,
    slnode/.style={draw, rounded corners, align=center, fill=blue!4,
                   inner sep=4pt, font=\footnotesize, minimum height=7mm},
    slplane/.style={draw, align=center, fill=green!5, inner sep=4pt,
                    font=\footnotesize, minimum height=7mm},
    sldec/.style={draw, diamond, aspect=2, align=center, fill=orange!8,
                  inner sep=1pt, font=\scriptsize},
    slflow/.style={-{Latex[length=2mm]}, thick},
    slbus/.style={-{Latex[length=2mm]}, thick, dashed},
  ]
    \node[slnode] (fc) {Forecast envelope\\(\texttt{smartload.forecast})};
    \node[sldec, below=of fc] (stale) {forecast\\stale?};
    \node[slplane, right=24mm of stale] (react) {Reactive fallback:\\observed RPS\\from telemetry};
    \node[sldec, below=of stale] (cmp) {compare $\hat{r}$\\to band $[(n{-}1)c,\,nc]$};
    \node[slnode, below left=10mm and 2mm of cmp] (out) {$\hat{r} > nc$\\\texttt{scale\_out}};
    \node[slnode, below=of cmp] (noop) {within band\\\texttt{noop}};
    \node[slnode, below right=10mm and 2mm of cmp] (in) {$\hat{r} < (n{-}1)c$\\\texttt{scale\_in}};
    \node[sldec, below=14mm of noop] (guard) {cooldown ok?\\within min/max?};
    \node[slplane, below=of guard] (act) {Docker SDK\\provision /\\decommission};
    \node[slnode, below=of act] (pub) {Publish \texttt{ScalingEvent}\\+ \texttt{scaling\_events} audit};

    \draw[slbus] (fc) -- (stale);
    \draw[slflow] (stale) -- node[above, font=\scriptsize] {yes} (react);
    \draw[slflow] (react) |- (cmp);
    \draw[slflow] (stale) -- node[right, font=\scriptsize] {no} (cmp);
    \draw[slflow] (cmp) -- (out);
    \draw[slflow] (cmp) -- (noop);
    \draw[slflow] (cmp) -- (in);
    \draw[slflow] (out) |- (guard);
    \draw[slflow] (in) |- (guard);
    \draw[slflow] (noop) -- (guard);
    \draw[slflow] (guard) -- node[right, font=\scriptsize] {pass} (act);
    \draw[slbus] (act) -- (pub);
  \end{tikzpicture}
  \caption{Forecast-driven autoscaler decision flow. A stale forecast diverts to
  a reactive observed-rate path; the dead band, cooldown, and min/max bounds
  suppress oscillation before the Docker SDK actuates and the event is audited.}
  \label{fig:imb:autoscaler}
\end{figure}

\section{Designing and Refining the Engines}\label{imb:iteration}

The engines did not arrive in their present form; each reached it through a
deliberate process of design, validation, and refinement, and the constant
through all of it was the deterministic baseline standing behind each trained
engine as a safety net. This closing section draws the three journeys together,
because the \emph{way} the engines were built is itself a result.

\paragraph{Anomaly detection: from a passing gate to a calibrated detector.} The
Isolation Forest engine cleared its $F_1 = 0.8012$ validation gate on the first
attempt. The work then went further and lined it up directly against the
threshold rule on production-shaped inputs through a comparison bench.
Calibrating the model and scaler into one shared coordinate system brought bench
agreement with the threshold baseline to $91.4\%$, with every remaining
disagreement in the value-adding direction. The threshold rule leads in
production for guaranteed stability and serves as the automatic fallback, while
the calibrated Isolation Forest is available as an enhanced detection mode. A
quorum guard in the sidecar protects both engines from an exclusion that would
empty the pool. Designing, validating, and then positioning each engine where it
is strongest is the same disciplined loop at work.

\paragraph{Forecasting: a dependable default with an advanced option in tuning.}
ARIMA$(3,0,1)$ is wired, serialised, and serving-ready, and stays under continued
tuning. The moving-average baseline remains the production default because it is
the reliable signal that drives the demonstrated autoscaling loop, and ARIMA sits
alongside it as the optional advanced predictor. The forecasting path works
end-to-end; the plan for the advanced predictor is concrete, namely a longer
training window and a hyperparameter search to close the MAPE gap.

\paragraph{Routing: a learned policy that matches a strong baseline.} The PPO
engine's serving path, safety gates, and benchmark harness are complete and
verified, and in shadow evaluation the policy ties round-robin, a strong,
industry-standard baseline. The training setup is well understood, and the next
direction is the binding constraint named in \cref{imb:ppo_finding}: retrain on a
heterogeneous-latency distribution, ideally with a closed-loop simulator, to give
the policy room to discriminate on varied traffic. The engine ships safely in
shadow in the meantime, with the deterministic fallback always available, and the
fuller analysis of the round-robin tie appears in \cref{ch:discussion}.

The thread running through all three is the same architectural decision made in
\cref{imb:pluggable}: because every learned engine is one swappable plugin behind
a deterministic baseline, SmartLoad can ship trained models confidently while
keeping production behaviour conservative. Whichever engine leads in production,
the baseline behind it is the same code path that has run safely from the first
sprint, and changing the lead is one line of configuration. That reversibility,
rather than any single model's score, is the implementation result this chapter
most wants to record.
